\section{Qualitative Analysis}
\label{sec:qualitative}

Figures~\ref{fig:qualitative-soldier}--\ref{fig:qualitative-3d} present
qualitative comparisons showing the geometric improvement produced by the
post-processor.

\begin{figure}[t]
  \centering
  \includegraphics[width=0.98\textwidth]{soldier_vis.png}
  \caption{%
    \textit{Soldier} sequence, r2 (0.05~bpp).
    \textbf{Left}: original (79,529 pts).
    \textbf{Centre}: PCGCv2 decoded (76,058 pts): surface protrusions and
    structural boundaries are smoothed away.
    \textbf{Right}: proposed post-processed output (76,058 pts): fine
    geometric structure substantially recovered at helmet, shoulder hardware,
    and facial boundaries.
    Point count is preserved by construction.
  }
  \label{fig:qualitative-soldier}
\end{figure}

\begin{figure}[t]
  \centering
  \includegraphics[width=0.98\textwidth]{soldier_3d_nocurv_ep24_r2.png}
  \caption{%
    \textit{Soldier} sequence 3D view, r2 (0.05~bpp).
    \textbf{Left}: original.
    \textbf{Centre}: PCGCv2 decoded: surface geometry is noticeably blurred
    at fine structures.
    \textbf{Right}: proposed: surface sharpness partially recovered.
  }
  \label{fig:qualitative-3d}
\end{figure}

\begin{figure}[t]
  \centering
  \includegraphics[width=0.98\textwidth]{redandblack_3d_nocurv_ep24_az10_r2.png}
  \caption{%
    \textit{Redandblack} sequence, 3D view at r2 (0.05~bpp).
    \textbf{Left}: original.
    \textbf{Centre}: PCGCv2 decoded: clothing texture and body contours
    are smoothed.
    \textbf{Right}: proposed: detail partially recovered across the surface.
  }
  \label{fig:qualitative-redandblack}
\end{figure}

The qualitative improvement is concentrated at geometric boundaries and surface
structures (helmet visor, shoulder hardware, clothing texture).
The overall body shape is unchanged; the model applies corrections at
high-curvature regions rather than globally rescaling the point cloud.

% ============================================================
